\section{From noisy conditional samples to the target law}
\label{sec:sampling}

The auxiliary algorithm returns a conditional sample with an added
Gaussian. We now use it to sample the original, unnoised law
$\pi_V$. The construction has two stages. A contracting proximal chain
first approaches $\widetilde\pi=\pi_V*N(0,hI_d)$. Conditional on its
output, a sequence of variance-halving steps reduces the remaining
sampling problem to a law close to a mode-centered Gaussian. The
numbers of calls in both stages are prescribed in advance.

Throughout this section, $V\in C^2(\R^d)$ satisfies
$V(0)=0$, $\nabla V(0)=0$, and
$I_d\preceq\nabla^2V\preceq2I_d$ everywhere. The origin is therefore
the supplied minimizer. Each original query returns both $V$ and
$\nabla V$ at the queried point. The main construction uses only the
gradient; the bounded-dimensional fallback also uses the value.
We fix $0<e<1$, put $\delta=d^{-e}$ and $h=\delta/2$, and take an
integer $K_0>2+1/e$. Choose $t>0$ as in \eqref{eq:t-choice}, and
fix the approximation orders required by Theorem~\ref{thm:primitive}.
All sufficiently-large-dimension thresholds below are uniform in $V$.

\subsection{A contracting noised chain}

The Gaussian proximal sampler alternates a Gaussian step with a
restricted Gaussian conditional draw
\citep[Definition~1 and Lemmas~2 and~4]{lee2021proximal}; see also
\citet[Section~3 and Theorem~1]{chen2022proximal}. We use the chain
on the noised variable and prove its contraction directly.

For $y\in\R^d$, denote by $\kappa_y$ the probability law with density
proportional to
\begin{equation}\label{eq:proximal-conditional}
              \exp\!\left(-V(x')-\frac{|x'-y|^2}{2h}\right).
\end{equation}
The ideal transition $P_h$ draws $x'\sim\kappa_y$ and returns
$x'+\sqrt h\,N$, where $N$ is an independent standard Gaussian.
Thus $P_h(y,\cdot)=\kappa_y*N(0,hI_d)$.

\begin{lemma}[Proximal contraction]\label{lem:proximal-contraction}
The law $\widetilde\pi$ is invariant for $P_h$, and
\[
 \Wtwo(\kappa_{y_1},\kappa_{y_2})
       \le\frac{|y_1-y_2|}{1+h}.
\]
Consequently, for input laws $\mu,\eta$ of finite second moment,
\[
 \Wtwo(\mu P_h,\eta P_h)
       \le\frac{\Wtwo(\mu,\eta)}{1+h}.
\]
\end{lemma}
\begin{proof}
Under the joint law $\pi_V(dx)N(x,hI_d)(dy)$, the conditional law
of $x$ given $y$ is $\kappa_y$. Drawing this conditional and then
redrawing $y$ given $x$ preserves the joint law, whose $y$-marginal
is $\widetilde\pi$.

The potential in \eqref{eq:proximal-conditional} has curvature at
least $1+1/h$. Replacing $y_1$ by $y_2$ changes its gradient by the
constant vector $(y_1-y_2)/h$. The synchronous comparison
\eqref{eq:synchronous-comparison} therefore bounds the conditional
$\Wtwo$ distance by $|y_1-y_2|/(1+h)$. Couple the subsequent Gaussian
increments identically. Integrating these conditional couplings
against an input coupling, as in Lemma~\ref{lem:coupling-smoothing},
proves the kernel contraction.
\end{proof}

To implement this transition, substitute $x'=y+\sqrt h\,z$ in
\eqref{eq:proximal-conditional}. The normalized conditional potential
becomes $|z|^2/2+F_y(z)$, where
\begin{equation}\label{eq:proximal-primitive}
 F_y(z)=V(y+\sqrt h\,z),\qquad
 g_y(z)=\nabla F_y(z)=\sqrt h\,\nabla V(y+\sqrt h\,z).
\end{equation}
The Hessian bound and $\nabla V(0)=0$ give
\[
 \Lip(g_y)\le2h=\delta,\qquad |g_y(0)|\le2\sqrt h\,|y|.
\]
An evaluation of $g_y$ is one original first-order query. Apply
Theorem~\ref{thm:primitive} to $F_y$ at center zero and scale one,
with fresh standalone seeds, and map its output $S$ to
$y+\sqrt h\,S$. This defines the actual transition $\widehat P_h$.
The primitive already approximates $Z+N$, so the ideal comparison
law of this affine image is exactly $P_h(y,\cdot)$.
In particular, the primitive's independent final reserve $G/2$
is retained. Its pre-reserve $\Wtwo$ guarantee and Gaussian smoothing
are what supply the following TV guarantee:
\begin{equation}\label{eq:proximal-local-errors}
 \begin{aligned}
 \Wtwo(\widehat P_h(y,\cdot),P_h(y,\cdot))
   &\le\sqrt h\,B_d\delta^{K_0}
                    (\sqrt d+2\sqrt h\,|y|),\\
 \norm{\widehat P_h(y,\cdot)-P_h(y,\cdot)}_{\TV}
   &\le B_d\delta^{K_0}(\sqrt d+2\sqrt h\,|y|).
 \end{aligned}
\end{equation}
Here and below $B_d=C(1+\log d)^C$ may be increased by a fixed
factor or power. Its constants depend only on the fixed parameters,
not on the center, oracle, or random history.

\begin{proposition}[Mixing of the approximate noised chain]
\label{prop:noised-mixing}
With $K_0>2+1/e$ and $t$ chosen above, start the
$\widehat P_h$-chain at zero and run
$\lceil Ch^{-1}\log d\rceil+1$ steps. For a sufficiently large
fixed $C$, its output has total-variation distance $o(1)$ from
$\widetilde\pi$, uniformly over all admissible $V$.
\end{proposition}
\begin{proof}
Lemma~\ref{lem:convex-preliminaries}, applied at the mode zero,
gives $\E_{\pi_V}|X|^2\le d$. Hence the root mean square under
$\widetilde\pi$ is at most $C\sqrt d$.
Let $\mu_k$ be the actual law after $k$ steps, and set
$a_k=\Wtwo(\mu_k,\widetilde\pi)$. Then
\[
 \left(\int |y|^2\,\mu_k(dy)\right)^{1/2}
       \le a_k+C\sqrt d.
\]
All these moments are finite. Indeed they are finite at $k=0$, and
the local $\Wtwo$ bound in \eqref{eq:proximal-local-errors}, together
with the conditional moment bounds for $P_h$, propagates finite
second moments to the next step.

Take the root mean square of the first local error in
\eqref{eq:proximal-local-errors}. Conditional coupling and
Lemma~\ref{lem:proximal-contraction} give
\[
 a_{k+1}\le
 \left(\frac1{1+h}+2hB_d\delta^{K_0}\right)a_k
       +C\sqrt h\,B_d\delta^{K_0}\sqrt d.
\]
For all large $d$, the coefficient of $a_k$ is at most $1-ch$
for an absolute $c>0$. Since $a_0\le C\sqrt d$, summing this
geometric recurrence yields
\begin{equation}\label{eq:proximal-global-error}
 a_k\le C\sqrt d\,e^{-chk}
       +Ch^{-1/2}B_d\delta^{K_0}\sqrt d.
\end{equation}
The second term is $o(\sqrt h)$: after division by $\sqrt h$,
it is bounded by a constant times
$B_d\delta^{K_0-1}\sqrt d$, which tends to zero because
$e(K_0-1)>1/2$. At $k=\lceil Ch^{-1}\log d\rceil$, a sufficiently
large fixed $C$ makes the first term $o(\sqrt h)$ as well.

Use one more transition to obtain TV control. The mixture bound
in Lemma~\ref{lem:coupling-smoothing} gives
\[
 \norm{\mu_k\widehat P_h-\mu_kP_h}_{\TV}
 \le B_d\delta^{K_0}
       \bigl[\sqrt d+2\sqrt h\,(a_k+C\sqrt d)\bigr]=o(1).
\]
For the remaining comparison, the conditional $x'$-laws from
$\mu_k$ and $\widetilde\pi$ can be coupled with root mean square
at most $a_k/(1+h)$. Their common independent $\sqrt h$-Gaussian
then gives, by \eqref{eq:gaussian-smoothing},
\[
 \norm{\mu_kP_h-\widetilde\pi P_h}_{\TV}
       \le\frac{a_k}{(1+h)\sqrt h}=o(1).
\]
Invariance of $\widetilde\pi$ completes the proof.
\end{proof}

\subsection{Conditional descent and its ideal histories}

An exact draw from $\kappa_y$ with $y\sim\widetilde\pi$ has
marginal $\pi_V$. We next implement this unnoised conditional draw.
For its analysis, first replace the incoming approximate $y$ by
$y\sim\widetilde\pi$. Applying the same conditional algorithm to
both inputs costs at most the TV error of
Proposition~\ref{prop:noised-mixing}, by data processing.

For a related recursive restricted-Gaussian-oracle framework, see
Chen--Chewi--Lu--Zhang \cite[Algorithm~3.3 and Section~6.3]{chen2026picard}.
The construction here uses fixed variance halving, mode-centered Gaussian
termination, and an every-execution query cap.

Fix $y$ temporarily. Keep $F_y$ and its oracle $g_y$ fixed throughout
the conditional algorithm; only the primitive center and scale
change. At center $x$ and scale $r$, the desired normalized law is
\[
 \nu_{x,r}(dz)\propto
             \exp\bigl(-|z|^2/2-F_y(x+rz)\bigr)\,dz.
\]
Suppose first that $Z\sim\nu_{x,r}$ and that $N$ is an independent
standard Gaussian. Given $Y=Z+N$, the conditional density of $Z$
is proportional to
\[
             \exp\bigl(-|z-Y/2|^2-F_y(x+rz)\bigr).
\]
Substitute $z=Y/2+z_{\rm new}/\sqrt2$ to obtain the exact identity
\begin{equation}\label{eq:conditional-descent}
 Z=\frac Y2+\frac{Z_{\rm new}}{\sqrt2},\qquad
 Z_{\rm new}\mid Y\sim\nu_{\,x+rY/2,\;r/\sqrt2}.
\end{equation}

The actual algorithm starts at $x_0=0,r_0=1$ and uses
$J=\lceil C_1\log d\rceil$ levels. At each $0\le j<J$, it calls
Theorem~\ref{thm:primitive} with fresh seeds at $(x_j,r_j)$ to
produce a noised observation $Y_j$, and sets
\[
 x_{j+1}=x_j+r_jY_j/2,\qquad r_{j+1}=r_j/\sqrt2.
\]
At level $J$, it produces a base sample $Z_J$ by the Gaussian
rule given below and returns
\[
                         y+\sqrt h\,(x_J+r_JZ_J).
\]
To identify this output, define the backwards variables
$Z_j=Y_j/2+Z_{j+1}/\sqrt2$ for the analysis. The center update gives
\[
 x_j+r_jZ_j
 =x_j+r_jY_j/2+(r_j/\sqrt2)Z_{j+1}
 =x_{j+1}+r_{j+1}Z_{j+1}.
\]
Since $x_0=0$ and $r_0=1$, the returned point is therefore
$y+\sqrt h\,Z_0$. By \eqref{eq:conditional-descent}, ideal
observations and an exact base conditional sample would make this
an exact draw from $\kappa_y$.

We must control the errors at random centers. The useful moment
bound concerns the history in which all preceding observations
have their ideal laws.

\begin{lemma}[Moments along the ideal descent]
\label{lem:descent-moments}
If $y\sim\widetilde\pi$ and every conditioning observation preceding
level $j$ is ideal, then
\[
 b_j:=\left(\E|g_y(x_j)|^2\right)^{1/2}\le C\sqrt d
             \qquad(0\le j\le J),
\]
with an absolute $C$ for all sufficiently large $d$.
\end{lemma}
\begin{proof}
The initial estimate follows from
$|g_y(0)|\le2\sqrt h\,|y|$ and the moment bound for
$\widetilde\pi$. Conditional on an ideal history, the next
observation has the law of $Z+N$ at center $x_j$ and scale
$r_j=2^{-j/2}$. Equation~\eqref{eq:primitive-moments} therefore gives
\[
 \left(\E[|Y_j|^2\mid x_j,y]\right)^{1/2}
          \le C(\sqrt d+r_j|g_y(x_j)|).
\]
The center update and $\Lip(g_y)\le\delta$ imply, by the
conditional and then unconditional $L^2$ triangle inequalities,
\[
 b_{j+1}\le(1+C\delta r_j^2)b_j+C\delta r_j\sqrt d.
\]
Both $\sum_{j\ge0}r_j$ and $\sum_{j\ge0}r_j^2$ are bounded
by absolute constants. In particular the products of
$1+C\delta r_j^2$ are bounded by
$\exp(C\delta\sum_jr_j^2)$. Expanding the recurrence proves
the stated uniform bound.
\end{proof}

Here is the precise replacement order used to sum the TV errors.
For $0\le j\le J$, let $H_j$ be the output law when the first
$j$ observations are ideal, all later observations use the actual
primitive, and the base rule is the actual Gaussian rule.
For $0\le j<J$, the laws $H_j$ and $H_{j+1}$ have the same ideal prefix of length
$j$. Conditional on that prefix, only the next observation kernel
changes; all subsequent computation is the same probability kernel.
Data processing and the mixture TV bound therefore give
\[
 \norm{H_j-H_{j+1}}_{\TV}
    \le B_d\delta^{K_0}\E_{\rm ideal}
                        \bigl[\sqrt d+|g_y(x_j)|\bigr].
\]
The expectation includes $y\sim\widetilde\pi$ and the common
ideal prefix. Lemma~\ref{lem:descent-moments} now yields
\begin{equation}\label{eq:descent-tv}
 \norm{H_0-H_J}_{\TV}
       \le CJ B_d\delta^{K_0}\sqrt d=o(1).
\end{equation}
Thus errors are summed using only ideal-history moments. It remains
to compare the actual and exact base laws under that same ideal history.

\subsection{The Gaussian base rule}

At depth $J$, write $l=\delta r_J^2$ and start at $z^{(0)}=0$.
Make exactly $m_0=\lceil C_2\log d\rceil$ gradient iterations
\begin{equation}\label{eq:mode-iteration}
 z^{(k+1)}=-r_Jg_y(x_J+r_Jz^{(k)}),
                         \qquad 0\le k<m_0.
\end{equation}
The base output is $Z_J=z^{(m_0)}+G$, for a fresh standard
Gaussian $G$. The iteration map has Lipschitz constant at most $l$.
Its unique fixed point is the mode $z_0$ of $\nu_{x_J,r_J}$, and
\eqref{eq:primitive-moments} gives
\begin{equation}\label{eq:mode-accuracy}
 |z^{(m_0)}-z_0|
       \le l^{m_0}|z_0|
       \le l^{m_0}\frac{r_J|g_y(x_J)|}{1-l}.
\end{equation}
The following comparison explains why a Gaussian centered there
completes the construction.

\begin{lemma}[Mode-centered Gaussian approximation]
\label{lem:mode-gaussian}
For any $F\in C^2(\R^d)$ with $\Lip(\nabla F)\le\lambda$, center
$x$, and scale $0<r\le1$ such that $l=\lambda r^2\le1/2$, let $z_0$ be
the mode of $\nu_{x,r}$. Then
\[
 \norm{\nu_{x,r}-N(z_0,I_d)}_{\TV}\le Cl\sqrt d.
\]
\end{lemma}
\begin{proof}
Translate by $z_0$. Up to an additive constant, the potential of
$w=z-z_0$ is $|w|^2/2+J(w)$, where
\[
 J(w)=z_0\cdot w+F(x+r(z_0+w))-F(x+rz_0).
\]
The mode equation gives $\nabla J(0)=0$, and
$\Lip(\nabla J)\le l$. In particular $J(0)=0$ and
$|J(w)|\le l|w|^2/2$.
For $0\le s\le1$, the potential $|w|^2/2+sJ(w)$ has mode zero
and Hessian between $(1-l)I_d$ and $(1+l)I_d$.
Its law $\mu_s$ satisfies, by Lemma~\ref{lem:convex-preliminaries},
\[
 \Var_{\mu_s}J
 \le\frac1{1-l}\E_{\mu_s}|\nabla J|^2
 \le\frac{l^2}{1-l}\E_{\mu_s}|w|^2
 \le\frac{l^2d}{(1-l)^2}\le4l^2d.
\]
The same quadratic bounds justify differentiating
\[
 \Phi(s)=\log\int_{\R^d}e^{-|w|^2/2-sJ(w)}\,dw
\]
twice under the integral, and give
$\Phi''(s)=\Var_{\mu_s}J$. Since $\mu_0=N(0,I_d)$,
\[
 \KL(\mu_0\Vert\mu_1)
  =\Phi(1)-\Phi(0)-\Phi'(0)
  =\int_0^1(1-s)\Phi''(s)\,ds\le2l^2d.
\]
Pinsker's inequality, followed by translation back by $z_0$,
proves the assertion.
\end{proof}

Under an ideal history, the Gaussian shift bound and
\eqref{eq:mode-accuracy} bound the expected TV error of the base rule
by
\begin{equation}\label{eq:descent-base-error}
 C\delta 2^{-J}\sqrt d
 +C(\delta2^{-J})^{m_0}
                 \frac{2^{-J/2}\sqrt d}{1-\delta2^{-J}}.
\end{equation}
Here the second term uses Lemma~\ref{lem:descent-moments}.
Choose the fixed $C_1$ large enough that
$\delta2^{-J}\sqrt d\to0$. A sufficiently large fixed $C_2$
makes the second term tend to zero as well. Replacing the base law
under the ideal history therefore changes TV by $o(1)$.
By \eqref{eq:conditional-descent}, the resulting completely ideal
conditional algorithm returns $\kappa_y$. Averaging over
$y\sim\widetilde\pi$ gives $\pi_V$.

Combining this base comparison, \eqref{eq:descent-tv}, and the
initial replacement of the incoming $y$ proves
\begin{equation}\label{eq:unnoised-output-error}
            \norm{\mathcal L(X_{\rm out})-\pi_V}_{\TV}=o(1),
\end{equation}
uniformly over admissible $V$. The final Gaussian is used to
approximate the last conditional law; the law in
\eqref{eq:unnoised-output-error} is the original target without
an additional convolution.

\subsection{Query counts and all dimensions}

\begin{theorem}[Subpolynomial upper bound]\label{thm:upper}
For every fixed $\varepsilon>0$, there is $C_\varepsilon<\infty$
such that, for every integer $d\ge2$, a measurable algorithm using
the exact value-and-gradient oracle makes at most
$C_\varepsilon d^\varepsilon$ queries on every execution and
returns a sample within TV distance $1/10$ of $\pi_V$ for every
$V$ satisfying \eqref{eq:original-class}. Thus
\[
                  Q(d)\le C_\varepsilon d^\varepsilon
                       \qquad(d\ge2).
\]
\end{theorem}
\begin{proof}
Given $\varepsilon>0$, first choose $e\in(0,1)$ with
$2e<\varepsilon$, then an integer $K_0>2+1/e$, and then
$t>0$ satisfying \eqref{eq:t-choice}. Fix the auxiliary construction's
orders in its prescribed order, and finally the constants in the
chain, descent depth, and base iteration count. These choices may
depend on $\varepsilon$, but they are fixed before observing the
oracle or the random seeds.

Every invocation of Theorem~\ref{thm:primitive} uses at most
$\polylog(d)\delta^{-t}$ evaluations of the gradient oracle $g_y$. By
\eqref{eq:proximal-primitive}, each evaluation is one original
query, including calls at random descent centers. The cap is uniform
in those centers. There are $O(\delta^{-1}\log d)$ proximal calls,
$O(\log d)$ descent calls, and $O(\log d)$ additional gradient
queries in the base rule. Summing gives the deterministic cap
\begin{equation}\label{eq:upper-query-count}
 \polylog(d)\delta^{-(1+t)}=d^{e(1+t)+o(1)}.
\end{equation}
The query counts and recursion shapes of this construction are
fixed by the scalar parameters; they do not depend on oracle
magnitudes or approximation events. All maps are measurable by
the primitive theorem and the explicit finite operations above.
Since $t<1$, we have $e(1+t)<2e<\varepsilon$.
The fixed polylogarithmic factors fit $C_\varepsilon d^\varepsilon$
for all sufficiently large $d$. Equation~\eqref{eq:unnoised-output-error}
gives TV error at most $1/10$ above a uniform dimension threshold.

For each of the finitely many smaller dimensions, use independent
standard Gaussian proposals. Accept a proposal $X$ with probability
\[
                    \exp\bigl(-V(X)+|X|^2/2\bigr).
\]
Integrating the Hessian bounds from the known minimizer gives
$|x|^2/2\le V(x)\le|x|^2$. The acceptance probability is therefore
at most one, and its average is at least
$\E e^{-|X|^2/2}=2^{-d/2}$. The proposal density times its
acceptance probability is proportional to $e^{-V(x)}$, so an
accepted proposal has law $\pi_V$.

Cap the number of trials at $\lceil4\cdot2^{d/2}\rceil$.
Return the first accepted proposal, or zero if none is accepted.
The failure probability is at most $e^{-4}<1/10$; conditional on
success, the returned proposal has law $\pi_V$. The fallback thus
has TV error at most $e^{-4}$ and makes at most the prescribed
number of original first-order queries on every execution, one per
trial. Increasing $C_\varepsilon$ covers these finitely many caps.
This proves the assertion for all $d\ge2$.
\end{proof}

Theorem~\ref{thm:upper} proves the upper assertion of
Theorem~\ref{thm:main}. Its approximation orders depend on the fixed
exponent $\varepsilon$; the result concerns oracle calls and imposes
no bound on the real computation between them. The next section
proves the independent logarithmic lower bound.
